\chapter{Dataset and Preprocessing}
\section{Raw Count Matrix}
The raw expression matrix contains genes as molecular features and cell barcodes as observations. Let $x_{ij}$ denote the raw count for gene $j$ in cell $i$. The resulting matrix is sparse because many transcripts are not detected in each cell.

\section{Quality Control}
Cells and genes were filtered to remove low-quality observations. Typical QC metrics included total counts per cell, number of detected genes per cell, and mitochondrial percentage. The mitochondrial percentage for cell $i$ is
\begin{equation}
MT_i = \frac{\sum_{j \in \mathcal{M}} x_{ij}}{\sum_{j=1}^{p} x_{ij}} \times 100,
\end{equation}
where $\mathcal{M}$ is the set of mitochondrial genes and $p$ is the number of genes.
\begin{table}[H]
\centering
\caption{QC filtering statistics template. Replace placeholder values with final notebook outputs if needed.}
\label{tab:qc_statistics}
\begin{tabular}{lrr}
\toprule
QC stage & Cells retained & Genes retained \\
\midrule
Raw matrix & -- & -- \\
After minimum gene filtering & -- & -- \\
After minimum cell filtering & -- & -- \\
After mitochondrial filtering & -- & -- \\
Final HVG matrix & -- & 3000 \\
\bottomrule
\end{tabular}
\end{table}


\begin{figure}[H]
\centering
\includegraphics[width=0.92\textwidth]{qc_workflow.pdf}
\caption{QC workflow showing filtering of low-quality cells, lowly detected genes, and high mitochondrial-percentage cells. The diagram is generated by \texttt{thesis/scripts/generate_thesis_figures.py}.}
\label{fig:qc_workflow}
\end{figure}

\section{Normalization and Log Transformation}
Library-size normalization was applied to reduce technical variation caused by differences in sequencing depth. For cell $i$, normalized expression is
\begin{equation}
\tilde{x}_{ij} = \frac{x_{ij}}{\sum_{k=1}^{p} x_{ik}} \times s,
\end{equation}
where $s=10,000$ is the target library size. Log transformation is then applied as
\begin{equation}
y_{ij}=\log(1+\tilde{x}_{ij}).
\end{equation}

\begin{figure}[H]
\centering
\includegraphics[width=0.9\textwidth]{normalization_workflow.pdf}
\caption{Normalization workflow from raw counts to library-size normalized and log-transformed expression.}
\label{fig:normalization_workflow}
\end{figure}

\section{Highly Variable Gene Selection}
Highly variable genes were selected to reduce dimensionality while retaining biologically informative variation. For gene $j$, the mean and variance are
\begin{equation}
\mu_j = \frac{1}{n}\sum_{i=1}^{n} y_{ij}, \qquad
\sigma_j^2 = \frac{1}{n-1}\sum_{i=1}^{n}(y_{ij}-\mu_j)^2.
\end{equation}
The top 3000 genes ranked by variance were retained for downstream modeling.

\begin{figure}[H]
\centering
\includegraphics[width=0.78\textwidth]{hvg_selection_plot.pdf}
\caption{Top highly variable genes selected by variance. Genes with higher variance are prioritized as candidate machine learning features.}
\label{fig:hvg_selection}
\end{figure}

\begin{figure}[H]
\centering
\includegraphics[width=0.78\textwidth]{mean_variance_plot.pdf}
\caption{Mean versus variance plot used to evaluate highly variable gene behavior across cells.}
\label{fig:mean_variance}
\end{figure}

\section{Dimensionality Reduction and UMAP}
UMAP visualization was used to inspect whether disease states show separable transcriptomic patterns. UMAP is not used as the primary classifier but provides exploratory evidence about global structure in the processed feature space.

\begin{figure}[H]
\centering
\includegraphics[width=0.82\textwidth]{umap_condition.pdf}
\caption{UMAP visualization colored by \texttt{condition.l1}. Separation between disease conditions suggests condition-associated transcriptomic structure.}
\label{fig:umap_condition}
\end{figure}
